\documentclass[12pt,a4paper]{report}

% --- PACKAGES ---
\usepackage[french]{babel}
\usepackage[T1]{fontenc}
\usepackage[utf8]{inputenc}
\usepackage{listings}
\usepackage{xcolor}
\usepackage[most]{tcolorbox}
\usepackage{geometry}

\geometry{margin=2.5cm}

% --- COULEURS ---
\definecolor{menu-blue}{rgb}{0.1, 0.2, 0.4}
\definecolor{code-bg}{rgb}{0.98, 0.98, 0.98}

% --- STYLE DES BOITES (Design corrigé pour Yahya) ---
\newtcolorbox{menubox}[1]{
    enhanced,
    attach boxed title to top left={yshift=-2mm, xshift=2mm},
    colback=white,
    colframe=menu-blue,
    fonttitle=\bfseries,
    colbacktitle=menu-blue,
    title=#1,
    boxed title style={size=small, sharp corners},
    drop shadow,
    bottom=2mm
}

% --- CONFIGURATION C ---
\lstset{
    language=C,
    backgroundcolor=\color{code-bg},
    basicstyle=\ttfamily\footnotesize,
    keywordstyle=\color{blue},
    commentstyle=\color{gray},
    stringstyle=\color{purple},
    breaklines=true,
    frame=none,
    numbers=left,
    numberstyle=\tiny\color{gray}
}

\begin{document}

\chapter{Module Système de Menus (Menus \& Actions)}

\section{Introduction}
L'architecture de navigation est le pilier de l'expérience utilisateur. Dans notre application \textit{Smart Transport Scheduler}, ce module gère la création dynamique de barres de menus et de panneaux contextuels. Il repose sur le concept de \textbf{GActionMap} de la GLib pour lier les clics de menu à des fonctions C.

\section{Architecture des Données (Header)}

\begin{menubox}{Structure : menu\_item\_config}
Pour gérer la récursivité (des menus dans des menus), la structure contient un pointeur vers ses propres enfants. L'énumération \texttt{kind} permet de distinguer une action simple, d'un sous-menu ou d'une section.
\end{menubox}

\begin{lstlisting}
typedef struct menu_item_config menu_item_config;

struct menu_item_config {
    const char *label;
    menu_item_kind kind;
    const menu_item_config *children;
    size_t child_count;
    bool enabled;
    menu_item_callback on_activate;
    gpointer user_data;
};
\end{lstlisting}

\section{Logique de Construction Récursive (Source)}

\begin{menubox}{Algorithme : append\_item\_recursive}
C'est la pièce maîtresse du module. Si un élément est de type \texttt{SUBMENU}, la fonction s'appelle elle-même pour construire le niveau inférieur. Cela permet de créer des menus d'une profondeur infinie de manière élégante.
\end{menubox}

\begin{lstlisting}
static void append_item_recursive(GMenu *menu, const menu_item_config *item, const menubar_config *config) {
    if (item == NULL || item->enabled == false) return;

    if (item->kind == MENU_ITEM_SUBMENU && item->children != NULL) {
        GMenu *submenu = g_menu_new();
        for (size_t i = 0; i < item->child_count; i++) {
            append_item_recursive(submenu, &item->children[i], config);
        }
        g_menu_append_submenu(menu, item->label, G_MENU_MODEL(submenu));
        g_object_unref(submenu);
        return;
    }
    // ... suite de la logique d'action ...
}
\end{lstlisting}

\section{Génération des Widgets}



\begin{menubox}{Fonction : create\_menubar}
Cette fonction transforme notre modèle logique en un widget \texttt{GtkPopoverMenuBar}. Elle sépare totalement la définition des menus (le "quoi") de leur affichage (le "comment"), respectant ainsi les principes du génie logiciel.
\end{menubox}

\begin{lstlisting}
GtkWidget *create_menubar(const menubar_config *config) {
    GMenuModel *model = build_menu_model(config);
    GtkWidget *bar = gtk_popover_menu_bar_new_from_model(model);
    
    if (config != NULL) {
        apply_widget_style(bar, &config->style);
    }
    
    g_object_unref(model);
    return bar;
}
\end{lstlisting}

\end{document}